% Part: applied-modal-logics
% Chapter: epistemic-logic
% Section: introduction

\documentclass[../../../include/open-logic-section]{subfiles}

\begin{document}

\olfileid[hi]{aml}{el}{int}

\olsection{परिचय}

जिस प्रकार मोडल तर्कशास्त्र \emph{मोडल प्रतिज्ञप्तियों} और उनके बीच के
निष्कर्षण-संबंधों का अध्ययन करता है, उसी प्रकार ज्ञानात्मक तर्कशास्त्र
\emph{ज्ञानात्मक प्रतिज्ञप्तियों} और उनके बीच के निष्कर्षण-संबंधों का अध्ययन
करता है। यहाँ एकस्थानी संयोजकों की व्याख्या संभावना और आवश्यकता को निरूपित
करने वाले मोडल संकारकों के रूप में नहीं, बल्कि ज्ञान या विश्वास का प्रतिरूपण
करने वाले ज्ञानात्मक अथवा विश्वासात्मक संकारकों के रूप में की जाती है।
उदाहरण के लिए, हम निम्न दावे व्यक्त करना चाह सकते हैं:

\begin{enumerate}
\item रिचर्ड जानता है कि कैलगरी अल्बर्टा में है।
\item ऑड्री मानती है कि संभवतः कोई कुत्ता सोफ़े पर है।
\item रिचर्ड जानता है कि ऑड्री जानती है कि उसकी कक्षा मंगलवार को होती है।
\item हर कोई जानता है कि एक वर्ष में 12 महीने होते हैं।
\end{enumerate}
समकालीन ज्ञानात्मक तर्कशास्त्र का आरंभ प्रायः याको हिंटिक्का की 1962 की
पुस्तक \emph{Knowledge and Belief} (ज्ञान और विश्वास) से माना जाता है। यह
उस समय लिखी गई थी जब तर्कशास्त्र में संभव-जगत अर्थविज्ञान का प्रयोग बढ़ रहा
था। वास्तव में ज्ञानात्मक तर्कशास्त्र अन्य मोडल तर्कशास्त्रों के अधिकांश
अर्थवैज्ञानिक साधनों का ही उपयोग करते हैं, किंतु उनकी व्याख्या अलग ढंग से
करते हैं। मुख्य परिवर्तन इस बात में है कि हम \emph{अभिगम्यता संबंध} को किसका
निरूपण मानते हैं। ज्ञानात्मक तर्कशास्त्रों में ये संबंध किसी प्रकार की
\emph{ज्ञानात्मक संभावना} को निरूपित करते हैं। हम देखेंगे कि जिस ज्ञानात्मक
धारणा का प्रतिरूपण किया जा रहा है, वह अभिगम्यता संबंध पर लगाई जाने वाली
शर्तों को प्रभावित करती है। हम यह भी देखेंगे कि अनुरूपता सिद्धांत को
ज्ञानात्मक व्याख्या देने पर क्या होता है। ऊपर के उदाहरणों में दो कर्ताओं,
रिचर्ड और ऑड्री, तथा दोनों के ज्ञान के बीच के संबंध का उल्लेख है। जिन
ज्ञानात्मक तर्कशास्त्रों पर हम विचार करेंगे वे बहु-कर्ता तर्कशास्त्र होंगे,
जिनमें ऐसी बातें व्यक्त की जा सकती हैं। इसके विपरीत, एकल-कर्ता ज्ञानात्मक
तर्कशास्त्र केवल इस बारे में बात करेगा कि कोई एक व्यक्ति क्या जानता या
मानता है।

\end{document}
